\begin{anexohorizontal}
\setlength{\textwidth}{\linewidth}
\begin{tablalafrox}{Reservas de contingencia y de gestión. Fuente: elaboración propia.}{tab_T15_contingencia}
  {L{2cm} L{2.2cm} L{3cm} L{2.2cm} Y}
  {\cab{Período} & \cab{HH de contingencia} & \cab{Peak adicional (personas equivalentes)} & \cab{Rol que más aporta} & \cab{Capacidad que la cubre}}
  Meses 1–6 & 5.000,78 & 11,6 en el mes 6 & DES, 6,7 & Desarrollo: 28 + 6,7 = 34,7, dentro de las 40 programadas. ARQ y DAT suman 17,4 equivalentes en el mes 5 frente a 15: ese exceso se nivela al mes 6, antes de la fecha límite del H3 \\
  Meses 7–12 & 5.253,15 & 11,8 en el mes 7 & DES, 8,6 & Desarrollo: 28 + 8,6 = 36,6, dentro de las 40 programadas \\
  Meses 13–15 & 3.230,04 & 21,2 en el mes 15 & IMP, 11,9 & Implantación: 23 + 11,9 = 34,9 frente a 30, con la opción de hasta 8 personas del contrato de implantación (sección 5.7) \\
  Meses 16–21 & 2.995,70 & 11,9 en el mes 20 & IMP, 11,8 & Implantación: 25 + 11,8 = 36,8 con la misma opción; calidad: 10 + 3,7 = 13,7 dentro de las 16 del mes 16 \\
  Meses 22–33 & 1.278,99 & 0,9 & SRE y SEG & SRE y NOC: dentro de las 44 personas \\
  Meses 34–56 & 2.440,32 & 0,9 & SRE y SEG & SRE y NOC: dentro de las 44 personas \\
\end{tablalafrox}
\end{anexohorizontal}
